\chapter[1999 --- Blobel]{1999: G\"unter Blobel}
\label{ch:1999_Blobel}

\section*{Main Contribution}

\textit{Discovered that proteins have intrinsic signal sequences that direct their transport and localization within the cell, explaining how cells organize their internal machinery.}\cite{nobel-1999,lecture-1999}

\section{Official Citation}

for the discovery that proteins have intrinsic signals that govern their transport and localisation in the cell

\section{Background and Historical Context}

A eukaryotic cell is a city of compartments. The nucleus houses the genome, the mitochondria generate energy, the endoplasmic reticulum synthesizes proteins and lipids, the Golgi apparatus processes and sorts these products, and lysosomes degrade cellular waste. Each compartment is enclosed by a membrane and contains a distinct set of proteins and metabolic activities. How do proteins, synthesized on ribosomes in the cytoplasm, find their way to the correct compartment? This question, which lies at the heart of cell biology, was answered by G\"unter Blobel, whose discovery of intrinsic protein sorting signals earned him the 1999 Nobel Prize in Physiology or Medicine.

G\"unter Blobel was born on May 21, 1936, in Waltersdorf, Silesia, Germany (now part of Poland). His family fled to West Germany in 1945 at the end of World War II, and Blobel grew up in Freiburg im Breisgau. He studied medicine at the University of T{\"u}bingen and then earned his Ph.D.\ in biochemistry at the University of Wisconsin in 1967 under the supervision of Mahlon Bush Kinter, working on lipid metabolism. In 1969, Blobel joined the laboratory of David Sabatini at the Rockefeller University in New York, where he began the work that would lead to his Nobel Prize--winning discovery.

The intellectual context of Blobel's work was shaped by the cell biology revolution of the 1960s and 1970s. George Palade and his colleagues had established the secretory pathway as the route by which proteins destined for the cell surface or for secretion are synthesized on the endoplasmic reticulum (ER), transported to the Golgi apparatus, and then delivered to their final destination by vesicular transport. Palade shared the Nobel Prize in 1974 for this work. However, the molecular mechanism by which proteins were directed to the ER in the first place---and how they were subsequently sorted within the cell---remained unknown.

\section{The Discovery/Contribution}

\subsection{The Signal Hypothesis}

Blobel's key insight, formulated in the early 1970s and published in a landmark paper in 1975, was the ``signal hypothesis.'' He proposed that proteins destined for secretion or for incorporation into certain organelles contain a short amino acid sequence---a ``signal sequence''---near their N-terminus, which directs the ribosome to the ER membrane and initiates translocation of the growing polypeptide chain across the ER membrane. The signal sequence, Blobel suggested, was recognized by a receptor on the ER membrane, which in turn recruited the ribosome and facilitated the passage of the polypeptide through a protein-conducting channel in the membrane.

The signal hypothesis was bold in its simplicity and far-reaching in its implications. It proposed that the information for protein localization was encoded within the protein itself---in its amino acid sequence---rather than being provided by external factors. This was a fundamentally new way of thinking about intracellular protein trafficking.

\subsection{Experimental Validation}

Blobel tested the signal hypothesis using an in vitro translation system---a cell-free system containing ribosomes, tRNAs, mRNAs, and other components necessary for protein synthesis. He showed that when mRNA encoding a secretory protein was translated in this system in the presence of microsomal membranes (small vesicles derived from the ER), the protein was translocated into the interior of the microsomes, just as it would be translocated into the ER in an intact cell. Moreover, the signal sequence was cleaved from the protein during translocation, as predicted by the signal hypothesis.

Blobel and his colleagues subsequently identified the molecular machinery responsible for protein translocation across the ER membrane. They discovered the signal recognition particle (SRP), a ribonucleoprotein complex that recognizes the signal sequence on the nascent polypeptide and targets the ribosome--nascent chain complex to the ER membrane. They also identified the SRP receptor, which mediates the docking of the ribosome--SRP complex to the ER membrane, and the protein translocation channel (the Sec61 complex), through which the polypeptide chain passes during translocation.

\subsection{Additional Sorting Signals}

As Blobel's work progressed, he and his colleagues discovered that proteins contain a variety of sorting signals that direct them to different cellular compartments. Proteins destined for the mitochondria contain a mitochondrial targeting signal (MTS), a positively charged amphipathic $\alpha$-helix at the N-terminus that directs the protein to the mitochondrial matrix. Proteins destined for the nucleus contain a nuclear localization signal (NLS), a short stretch of basic amino acids that directs the protein through nuclear pores. Proteins destined for peroxisomes contain a peroxisomal targeting signal (PTS).

Blobel also discovered the process of membrane insertion---the mechanism by which proteins become embedded in the lipid bilayer of cellular membranes. He showed that transmembrane proteins contain hydrophobic transmembrane segments that serve as signal-anchor sequences, directing the ribosome to the ER membrane and initiating the insertion of the protein into the membrane.

\subsection{The Unifying Principle}

The overarching principle that emerged from Blobel's work is that the intracellular localization of a protein is determined by intrinsic sorting signals encoded in its amino acid sequence. This principle unifies our understanding of protein trafficking across all cellular compartments and provides the molecular framework for the secretory pathway, the endomembrane system, and the targeting of proteins to mitochondria, chloroplasts, peroxisomes, and the nucleus.

\section{Impact on Modern Medicine}

Blobel's discoveries have had a significant impact on cell biology, molecular medicine, and biotechnology.

\subsection{Genetic Disease}

Mutations that disrupt protein sorting signals cause a range of genetic diseases. Cystic fibrosis, the most common lethal genetic disease in Caucasian populations, is caused by mutations in the CFTR gene that result in misfolding and mislocalization of the CFTR chloride channel protein. The protein is retained in the ER and degraded by the proteasome, rather than being transported to the cell surface. Understanding the ER quality control machinery---the system that recognizes and degrades misfolded proteins---has guided the development of CFTR modulators such as ivacaftor and lumacaftor, which correct the trafficking defect.

Similarly, mutations in sorting signals are responsible for certain forms of hypercholesterolemia (mutations in the LDL receptor that prevent its proper trafficking to the cell surface) and for some forms of lysosomal storage diseases (mutations in lysosomal enzyme sorting signals that prevent the enzymes from reaching the lysosome).

\subsection{Biotechnology and Pharmaceutical Production}

The understanding of protein sorting signals has been essential for the biotechnology industry. The production of therapeutic proteins---including insulin, growth hormone, clotting factors, and monoclonal antibodies---in mammalian cell culture systems depends on the ER secretory pathway. By fusing the gene encoding a therapeutic protein to a signal sequence, biotechnologists can direct the protein into the secretory pathway, where it is folded, modified, and secreted into the culture medium for purification.

\subsection{Understanding of Neurodegenerative Disease}

The ER quality control system that Blobel's work helped elucidate is implicated in many neurodegenerative diseases. In Huntington's disease, the mutant huntingtin protein aggregates in the ER and disrupts protein homeostasis. In Alzheimer's disease, the accumulation of misfolded amyloid-$\beta$ and tau proteins overwhelms the ER's folding capacity, triggering the unfolded protein response and contributing to neuronal death. Understanding these mechanisms is guiding the development of therapies that target the ER stress response.

\section{Legacy}

G\"unter Blobel spent his entire career at the Rockefeller University, where he became John Simon Guggenheim Memorial Foundation Professor and head of the Laboratory of Biochemistry and Cell Biology. He received numerous honors, including the Albert Lasker Basic Medical Research Award in 1993 and the Gairdner Foundation International Award in 1986. Blobel was also known for his commitment to science education and for his efforts to preserve the historic architecture of the Rockefeller University campus. He passed away on February 18, 2018, at the age of 81.

Blobel's discovery that proteins contain intrinsic sorting signals resolved one of the most fundamental questions in cell biology: how do proteins find their way to the correct address in the cell? The answer---that the address is written into the protein's own amino acid sequence---elegantly connects the information in the genome to the organization of the living cell. This principle continues to guide research in cell biology, medicine, and biotechnology and stands as a lasting legacy of one of the twentieth century's most influential cell biologists.


\section{Modern Developments}

Blobel's signal sequence work has led to modern protein trafficking research. Understanding of protein sorting has enabled development of targeted drug delivery systems. Understanding of nuclear pore complexes has informed development of nuclear-targeted therapies. Understanding of ER targeting has enabled production of recombinant proteins in bacteria and mammalian cells. Understanding of mitochondrial targeting has led to development of mitochondrial-targeted antioxidants.


\section{Bibliography}

\begin{thebibliography}{9}
\bibitem{nobel-1999}
Nobel Prize Outreach, ``The Nobel Prize in Physiology or Medicine 1999,''\emph{NobelPrize.org}.\\
\url{https://www.nobelprize.org/prizes/medicine/1999/summary/}

\bibitem{lecture-1999}
G"unter Blobel, ``Nobel Lecture,''\emph{NobelPrize.org}. Winner-authored primary source; downloadable from the official Nobel Prize archive.\\
\url{https://www.nobelprize.org/prizes/medicine/1999/blobel/lecture/}
\end{thebibliography}
